\section{\texttt{TCOLMAX}}
\noindent\textbf{Descriptor profile.} Vec entry 13, tile catalog opcode \texttt{0x1018}. Descriptor bundle: \texttt{B.OP + B.DIM + B.IOT}; WO/WM sites: 4; VEC beat-level sites: 4; ROM bytes: 32.\par\smallskip
\TileInstructionEncoding{figs/encoding/tile_instruction_entries/entry_13_pto_tcolmax_th0.pdf}{figs/encoding/tile_instruction_entries/entry_13_pto_tcolmax_th2.pdf}{0x1018}{13}{B.OP + B.DIM + B.IOT}
\paragraph{Bundled descriptor (PTO-ISA header).}
\begin{description}[leftmargin=0.18\linewidth,style=nextline]
\item[Bundle] \texttt{B.OP + B.DIM + B.IOT}. \texttt{B.OP.desc\_count} = 2; \texttt{B.OP.rom\_entry\_id} = 13.
\item[Assembly] 0: \texttt{B.OP} selects \texttt{TCOLMAX}, carries element format/variant, declares \texttt{desc\_count}, and names the ROM entry. 1: \texttt{B.DIM} supplies column-axis reduction or column-broadcast bounds. 2: \texttt{B.IOT} binds architectural tile operands. Across all \texttt{B.IOT} descriptors: tile inputs \texttt{src}, mask inputs none, and outputs \texttt{dst}.
\item[Tile operands] 1 tile input(s): \texttt{src}; 1 tile output(s): \texttt{dst}. Bound by 1 architectural \texttt{B.IOT} descriptor; each \texttt{B.IOT} supplies at most two source selectors plus one destination selector.
\item[Scalar GPR operands] No scalar GPR import/export descriptor is required.
\item[Metadata operands] \texttt{B.DIM} supplies column-axis reduction or column-broadcast bounds.
\end{description}
\par\smallskip
{\small
\paragraph{PTO ISA description.}
Reduce each column by taking the maximum across rows.
\paragraph{PTO ISA operation.}
Let \texttt{R = src.GetValidRow()} and \texttt{C = src.GetValidCol()}. For \texttt{0 \textless{}= j \textless{} C}:

\[
\mathrm{dst}_{0,j} = \max_{0 \le i < R} \mathrm{src}_{i,j}
\]
}\par\smallskip
\paragraph{Microcode site table.}
{\scriptsize
\begin{longtable}{@{}R{0.055\linewidth}L{0.23\linewidth}L{0.31\linewidth}L{0.22\linewidth}@{}}
\toprule
Seq & Micro instruction & WO[63:32] fields & WM[31:0] fields \\
\midrule
0 & \texttt{u\_vec.vlds} & \texttt{opc=0x17, x=0, fl=mem, seq=0, kind=b} & \texttt{entry=13, cnt=2, res=vec, var=0} \\
1 & \texttt{u\_vec.vlds} & \texttt{opc=0x17, x=0, fl=mem, seq=1, kind=b} & \texttt{entry=13, cnt=2, res=vec, var=0} \\
2 & \texttt{u\_vec.vmax} & \texttt{opc=0x1A, x=0, fl=alu, seq=2, kind=b} & \texttt{entry=13, cnt=2, res=vec, var=0} \\
3 & \texttt{u\_vec.vsts} & \texttt{opc=0x2A, x=0, fl=mem, seq=3, kind=b} & \texttt{entry=13, cnt=2, res=vec, var=0} \\
\bottomrule
\end{longtable}
}
\paragraph{ROM body DSL.}
\begin{lstlisting}[style=ptocode]
def rom_tcolmax(src, dst):
    dtype = dst.element_type
    valid_rows, valid_cols = src.valid_shape

    lanes = pto.get_lanes(dtype)
    remained = valid_cols

    for col_chunk in range(0, valid_cols, lanes):
        mask, remained = pto.make_mask(dtype, remained)

        acc = pto.vlds(src[0, col_chunk:])
        for row in range(1, valid_rows, 1):
            row_vec = pto.vlds(src[row, col_chunk:])
            acc = pto.vmax(acc, row_vec, mask)
        pto.vsts(acc, dst[0, col_chunk:], mask)

    return
\end{lstlisting}
\Needspace{0.34\textheight}
